% =====================================================================
%  CAPÍTULO 4 · ANÁLISIS Y DISEÑO
%  Función: definir QUÉ debe hacer el sistema (requisitos verificables) y
%  CÓMO se ha decidido construirlo (arquitectura y decisiones justificadas).
%  Reglas clave:
%   - Requisitos con ID (RF-01, RNF-01), atómicos y verificables. Di de
%     dónde salen (estado del arte, entrevistas, normativa...).
%   - Lo que no has medido no se afirma como cumplido («cumple WCAG AA»
%     solo si lo has comprobado y lo muestras).
%   - Cada diagrama se explica en el texto componente a componente.
%   - Trabajos con datos/ML: añade aquí (o en el Cap. 5) una ficha por
%     dataset: estructura, nº de muestras y particiones, tipo de dato con
%     un ejemplo visual, referencia y protocolo de evaluación.
% =====================================================================
\chapter{Análisis y diseño}
\label{cap:diseno}

Este capítulo especifica los requisitos de \emph{SalaLibre} y describe la
arquitectura y las decisiones de diseño que se tomaron para cumplirlos.

\section{Requisitos}
\label{sec:requisitos}

Los requisitos se obtuvieron de la comparativa del
Capítulo~\ref{cap:estado} y de una entrevista con el personal de la biblioteca.
Cada requisito tiene un identificador que se usa en las pruebas del
Capítulo~\ref{cap:pruebas}. La Tabla~\ref{tab:rf} recoge los requisitos
funcionales y la Tabla~\ref{tab:rnf}, los no funcionales.

\begin{table}[htbp]
  \centering
  \caption{Requisitos funcionales (extracto).}
  \label{tab:rf}
  \small
  \begin{tabularx}{\textwidth}{@{}lLl@{}}
    \toprule
    \textbf{ID} & \textbf{Descripción} & \textbf{Objetivo} \\
    \midrule
    RF-01 & El usuario puede consultar las salas libres en este momento. & \ref{obj:app} \\
    RF-02 & El usuario puede reservar una sala libre en dos pasos como máximo. & \ref{obj:app} \\
    RF-03 & El usuario recibe una notificación 10 minutos antes de su reserva. & \ref{obj:app} \\
    RF-04 & El sistema impide dos reservas solapadas de la misma sala. & \ref{obj:app} \\
    \bottomrule
  \end{tabularx}
  \fuente{Elaboración propia.}
\end{table}

\begin{table}[htbp]
  \centering
  \caption{Requisitos no funcionales (extracto).}
  \label{tab:rnf}
  \small
  \begin{tabularx}{\textwidth}{@{}lLl@{}}
    \toprule
    \textbf{ID} & \textbf{Descripción y criterio de aceptación} & \textbf{Objetivo} \\
    \midrule
    RNF-01 & Rendimiento: latencia media $<$ \qty{100}{ms} con 200 usuarios concurrentes. & \ref{obj:rendimiento} \\
    RNF-02 & Usabilidad: puntuación \ac{SUS} media $\geq 70$ con al menos 8 usuarios. & \ref{obj:validacion} \\
    RNF-03 & Seguridad: las contraseñas se almacenan con una función de \emph{hash} con sal. & \ref{obj:app} \\
    \bottomrule
  \end{tabularx}
  \fuente{Elaboración propia.}
\end{table}

% GUÍA: «El sistema será escalable» NO es verificable. Da un criterio medible
% o márcalo explícitamente como objetivo de diseño no medido.

\section{Arquitectura del sistema}
\label{sec:arquitectura}

La Figura~\ref{fig:arquitectura} muestra la arquitectura en tres capas. La capa
de presentación (arriba) es la aplicación Android, que solo se comunica con el
servidor mediante la \ac{API}. La capa de lógica (centro) valida las reservas y
envía las notificaciones. La capa de datos (abajo) guarda salas, usuarios y
reservas. Se eligió separar la lógica en el servidor para que la regla RF-04
(sin reservas solapadas) se compruebe en un único punto, dentro de una
transacción de la base de datos.

\begin{figure}[htbp]
  \centering
  \begin{tikzpicture}[
      capa/.style={draw, rounded corners, minimum width=11.8cm, minimum height=1.3cm,
                   align=center, fill=#1},
      comp/.style={draw, rounded corners, fill=white, minimum width=2.7cm,
                   minimum height=0.7cm, font=\small},
      flecha/.style={-{Stealth[length=2.5mm]}, thick}]
    \node[capa=RoyalBlue!10] (pres) {};
    \node[anchor=west, font=\small\bfseries] at (pres.west) {\ Presentación};
    \node[comp] at ([xshift=0.6cm]pres.center) (ui) {Pantallas (UI)};
    \node[comp, right=0.4cm of ui] (vm) {Estado de la vista};
    \node[capa=ForestGreen!10, below=0.6cm of pres] (log) {};
    \node[anchor=west, font=\small\bfseries] at (log.west) {\ Lógica};
    \node[comp] at ([xshift=0.6cm]log.center) (api) {API REST};
    \node[comp, right=0.4cm of api] (not) {Notificaciones};
    \node[capa=BrickRed!10, below=0.6cm of log] (dat) {};
    \node[anchor=west, font=\small\bfseries] at (dat.west) {\ Datos};
    \node[comp] at ([xshift=0.6cm]dat.center) (db) {Base de datos};
    \draw[flecha] (ui) -- (api);
    \draw[flecha] (api) -- (db);
    \draw[flecha] (not.north) -- (vm.south);
  \end{tikzpicture}
  \caption{Arquitectura en capas de \emph{SalaLibre}.}
  \label{fig:arquitectura}
  \fuente{Elaboración propia.}
\end{figure}

\section{Diseño de la interfaz}
\label{sec:diseno-interfaz}

En una aplicación, muestra los prototipos (\emph{wireframes}) y el mapa de
navegación, y relaciónalos con los requisitos. La
Figura~\ref{fig:navegacion} resume la navegación: desde la pantalla de inicio se
llega a la reserva en dos pasos, como exige RF-02.

\begin{figure}[htbp]
  \centering
  \begin{tikzpicture}[node distance=1.6cm,
      pant/.style={draw, rounded corners, fill=gray!8, minimum width=2.4cm,
                   minimum height=0.9cm, align=center, font=\small},
      flecha/.style={-{Stealth[length=2.5mm]}, thick}]
    \node[pant] (login) {Inicio de\\sesión};
    \node[pant, right=of login] (inicio) {Inicio\\(salas libres)};
    \node[pant, right=of inicio] (reserva) {Nueva\\reserva};
    \node[pant, below=0.8cm of inicio] (mis) {Mis\\reservas};
    \draw[flecha] (login) -- (inicio);
    \draw[flecha] (inicio) -- node[above, font=\scriptsize] {paso 1} (reserva);
    \draw[flecha] (reserva.south) |- node[pos=0.25, right, font=\scriptsize] {paso 2} (mis.east);
    \draw[flecha] (inicio) -- (mis);
  \end{tikzpicture}
  \caption{Mapa de navegación de la aplicación.}
  \label{fig:navegacion}
  \fuente{Elaboración propia.}
\end{figure}
